\documentclass[11pt]{article}

\usepackage[T1]{fontenc}
\usepackage{lmodern}
\usepackage[margin=1in]{geometry}
\usepackage{amsmath}
\usepackage{amssymb}
\usepackage{booktabs}
\usepackage{longtable}
\usepackage{array}
\usepackage{hyperref}

\hypersetup{
  colorlinks=true,
  linkcolor=blue!50!black,
  urlcolor=blue!60!black,
  citecolor=blue!50!black
}

\title{Text Property Graphs for Security Text}
\author{}
\date{}

\begin{document}
\maketitle

\section{Introduction}

Security text is not only a sequence of words. A CVE description, an
exploit report, or a vendor advisory states which vulnerability is being
discussed, which software is affected, what version is vulnerable, how
the attack works, what the impact is, and what mitigation is recommended.
These facts are not independent keywords. They are connected by syntax,
sentence order, entity references, discourse markers, and security
relations.

The Text Property Graph (TPG) was designed to preserve that structure.
It represents text as a directed, typed, attributed property graph.
Nodes represent linguistic or domain-specific objects: documents,
sentences, tokens, entities, predicates, CVE identifiers, software
names, attack vectors, impacts and remediations. Edges represent typed
relations between those nodes: containment, dependency, token order,
sentence order, coreference, predicate-argument binding, discourse
relation, entity relation, and optional security-specific relations.

The design is inspired by the Code Property Graph (CPG). A CPG combines
several views of source code (syntax, control flow, data flow) into one
graph. The TPG applies the same idea to natural-language text. Instead
of keeping dependency parsing, named entity recognition, coreference,
discourse structure and domain relations as separate outputs, the TPG
stores them as different edge types over one shared graph.

This chapter explains the problem that motivates TPG, the core
mechanics of the representation, the design choices we made along the
way (GNN over standard models, TPG over plain transformers, multiview
fusion over a single GNN), and the linguistic backend (spaCy) that
turns raw text into the initial nodes and edges. The claims here are
limited to what is implemented and documented in the project.

\section{Why We Chose a GNN-Based Solution for EPSS}

EPSS prediction is a binary or soft-label classification problem on a
single CVE: given everything we know about the CVE, predict the
probability that it will be exploited in the wild. The natural question
is: what kind of model fits this best?

A CVE record is not one thing. It is a small bundle of heterogeneous
inputs. There is a free-text description (one or two sentences). There
is a CVSS vector (eight categorical metrics). There is a CWE
classification (a list of weakness IDs). There are publication dates,
exploit-code flags, and source-mention counts. And, in our datasets,
there are LLM summary columns. Some of these inputs are tabular.
Some are sequential text. Some are graph-shaped relations between
entities mentioned in the text.

We chose a graph neural network because it is the model class that
handles heterogeneous, structurally varied input most directly:

\begin{itemize}
    \item A GNN's input is a graph. We can represent each CVE as a
    graph whose nodes carry features and whose edges carry typed
    relations. The shape and size of the graph can vary per CVE; the
    GNN does not require a fixed input length.
    \item A GNN learns from local neighbourhoods through message
    passing. This means the prediction for a CVE depends not only on
    which words appear, but on which entities are connected to which
    other entities. That is closer to how a security analyst reads an
    advisory.
    \item A GNN extends naturally to a hybrid model. We can attach
    contextual text embeddings to nodes and concatenate the pooled graph
    embedding with a tabular feature vector inside the same model. The
    graph branch sees the typed relational structure, the tabular branch
    sees the structured CVE metadata, and the head sees both.
\end{itemize}

A second reason is interpretability. With a GNN over a TPG, we can
trace which nodes received the highest attention weight and which edges
carried the most useful signal. With a flat sequence model the same
question becomes much harder.

A third reason is room to grow. If we later want to link a CVE to its
patch commit, to a CWE knowledge graph, or to source-code constructs
extracted by Joern, we can add new node types and edge types without
changing the rest of the pipeline.

\section{Why We Chose TPG to Drive the GNN}

Once we settled on a GNN, the next question was: what graph?

A few options were on the table. We could have used a dependency parse
alone. We could have used an Abstract Meaning Representation (AMR)
graph. We could have built a knowledge graph of typed entity triples.
Each of these is a single-view representation: a parse captures only
syntax, AMR captures only predicate-argument structure, a knowledge
graph captures only typed relations between named entities.

The Code Property Graph (CPG) introduced by Joern for source-code
analysis offered a different idea. A CPG keeps several views of a
program (Abstract Syntax Tree for syntax, Control Flow Graph for
execution order, Reaching-Definitions Graph for data flow, Control
Dependence Graph for guarded statements) as different edge types over
one shared node set. The benefit is that a query can cross views
without joining separate graphs.

We borrowed that pattern for text. The TPG keeps the four CPG-style
views as different edge groups over the same nodes:

\begin{table}[h]
\centering
\small
\begin{tabular}{ll}
\toprule
CPG view & Text-domain analogue used in the TPG \\
\midrule
Abstract Syntax Tree (AST) & Dependency parse, edge type \texttt{DEP} \\
Control Flow Graph (CFG)   & Token, sentence, paragraph order: \texttt{NEXT\_TOKEN}, \texttt{NEXT\_SENT}, \texttt{NEXT\_PARA} \\
Reaching-Definitions Graph & Coreference, edge type \texttt{COREF} \\
Control Dependence Graph   & Discourse and rhetorical relations: \texttt{RST\_RELATION}, \texttt{DISCOURSE} \\
\bottomrule
\end{tabular}
\caption{The four CPG views and their text-domain analogues in the TPG.}
\end{table}

To these we added structural edges (\texttt{CONTAINS}, \texttt{BELONGS\_TO})
and predicate-argument edges (\texttt{SRL\_ARG}, \texttt{ENTITY\_REL},
\texttt{SIMILARITY}, \texttt{AMR\_EDGE}) that are not in the CPG but
are useful for text. The full mapping is recorded in the schema as a
machine-readable dictionary so the correspondence is unambiguous.

The reason this design helps is the same reason it helped Joern. A
single query (or, in our case, a single GNN forward pass) can use
information from all views at once. The model is not forced to pick
between syntax, sequence, coreference and discourse. It sees them
together as different edges in the same graph.

\section{The TPG Is Generic, Not Just a Security Representation}

We use the TPG for security text in this project, but the schema
itself is not security-specific. The bottom layer is generic English
NLP. The 13 base node types (\textsc{document}, \textsc{paragraph},
\textsc{sentence}, \textsc{token}, \textsc{entity}, \textsc{predicate},
\textsc{argument}, \textsc{concept}, \textsc{noun\_phrase},
\textsc{verb\_phrase}, \textsc{clause}, \textsc{mention},
\textsc{topic}) and the 13 base edge types (\texttt{DEP},
\texttt{NEXT\_TOKEN}, \texttt{NEXT\_SENT}, \texttt{NEXT\_PARA},
\texttt{COREF}, \texttt{SRL\_ARG}, \texttt{AMR\_EDGE},
\texttt{RST\_RELATION}, \texttt{DISCOURSE}, \texttt{CONTAINS},
\texttt{BELONGS\_TO}, \texttt{ENTITY\_REL}, \texttt{SIMILARITY})
describe any English text. They make no assumption about the topic.

The security part is a Level 2 overlay built on top of that generic
layer. It adds 12 domain node types (CVE\_ID, CWE\_ID, SOFTWARE,
VERSION, COMPONENT, CODE\_ELEMENT, ATTACK\_VECTOR, IMPACT,
REMEDIATION, THREAT\_ACTOR, VULN\_TYPE, SEVERITY) and 10
\texttt{SEC\_*} edge types. The overlay is optional. When it is not
loaded, the schema is exactly the generic Level 1 schema and works on
any English text without changes.

A different domain can be added the same way:

\begin{itemize}
    \item For \textbf{clinical text}, an overlay would add node types
    such as \textsc{drug}, \textsc{dose}, \textsc{route},
    \textsc{condition}, \textsc{symptom} and edge types such as
    \texttt{TREATS}, \texttt{HAS\_DOSE}, \texttt{ADMINISTERED\_BY},
    \texttt{CAUSES\_SIDE\_EFFECT}.
    \item For \textbf{legal text}, an overlay would add
    \textsc{statute}, \textsc{case}, \textsc{party}, \textsc{holding}
    with edges \texttt{CITES}, \texttt{OVERRULES}, \texttt{DISTINGUISHES}.
    \item For \textbf{scientific text}, an overlay would add
    \textsc{method}, \textsc{dataset}, \textsc{result},
    \textsc{citation} with edges \texttt{EVALUATES\_ON},
    \texttt{COMPARES\_AGAINST}, \texttt{REPORTS\_RESULT}.
\end{itemize}

In every case the generic layer (sentences, tokens, dependency parse,
coreference, discourse) is shared. Only the typed entities and the
typed relations between them change. Anything built on top of the
generic layer (the GNN backbone, the multiview decomposition, the
hybrid head) keeps working without modification, because the model
consumes a typed property graph regardless of which types are present.

In other words, the TPG is a general-purpose representation for
information processing on text. The security extension is one
instantiation of it. The same machinery would carry medical case
notes, legal opinions, scientific abstracts, or any other domain whose
meaning depends on typed entities and typed relations between them.

\section{Why a GNN over a TPG, Not a Pure Transformer or a Pure Tabular MLP}

A reasonable alternative would have been to skip the graph and use a
standard NLP transformer (BERT or SecBERT) on the description plus a
tabular MLP on the CVSS one-hot and the count features, then
concatenate the two embeddings. Why did we go further?

\paragraph{Pure transformer on text alone.}
A transformer reads the description and produces a fixed-length
embedding. That embedding is dense and contextual, but it does not
expose any structure. If we want to know which token in the
description corresponds to the CVE identifier, the affected software,
or the attack vector, the transformer does not say. The information
is implicit in the attention weights but is not surfaced as queryable
structure. Adding the CVSS features as tabular input next to a
transformer embedding works (and is in fact what our model does in
the tabular branch), but it gives up the relational structure inside
the text. A typical CVE
description is short (one or two sentences) and dense with relations,
so losing that structure has a real cost.

\paragraph{Pure tabular MLP with one-hot encodings.}
We could also have skipped the text altogether and trained an MLP on
just the CVSS one-hot, the CWE multi-hot, and the count features. This
is the simplest and most leaky-resistant baseline. The risk is the
opposite of the transformer case: we drop everything the description
says about the vulnerability. Two CVEs with the same CVSS vector and
no CWE list look identical to a tabular MLP. The description is what
distinguishes them.

\paragraph{What we actually do.}
The TPG-plus-GNN-plus-tabular hybrid sits between these extremes. The
graph branch reads the description through the TPG, gets a 768-dim
SecBERT embedding per token (so the contextual signal is preserved on
the nodes), and runs message passing over typed edges. The tabular
branch reads the CVSS one-hot, the CWE multi-hot, age, and counts as a
55-dim vector. The two branches meet in the prediction head. Neither
the structure nor the metadata is thrown away.

The 20-run text-source ablation reports that on the in-distribution
corpus the description carries the predictive signal and the LLM
summaries do not add to it. That is an empirical finding about the
data, not an indictment of the architecture. The architecture supports
a clean reading of the result because the model produces interpretable
per-variant deltas.

\section{Multiview GNN, and Why ``Multiview''}

The GNN backbone we use is called the Multiview GNN.

\paragraph{What ``multiview'' means here.}
A view is a subset of the edge-type vocabulary. The model partitions
the 13 generic edge types (and the 10 \texttt{SEC\_*} edge types when
security edges are enabled) into five groups, called views:

\begin{table}[h]
\centering
\small
\begin{tabular}{ll}
\toprule
View & Edge types in the view \\
\midrule
Syntactic  & \texttt{DEP}, \texttt{CONTAINS}, \texttt{BELONGS\_TO} \\
Sequential & \texttt{NEXT\_TOKEN}, \texttt{NEXT\_SENT}, \texttt{NEXT\_PARA} \\
Semantic   & \texttt{COREF}, \texttt{SRL\_ARG}, \texttt{AMR\_EDGE} \\
Discourse  & \texttt{RST\_RELATION}, \texttt{DISCOURSE}, \texttt{ENTITY\_REL}, \texttt{SIMILARITY} \\
Security   & all \texttt{SEC\_*} types (only present when security edges are enabled) \\
\bottomrule
\end{tabular}
\caption{The five edge-type views used by the multiview GNN.}
\end{table}

For each view, the model runs a separate Gated Graph Convolution
(GGNN) restricted to that view's edges. So the syntactic view sees
only \texttt{DEP}, \texttt{CONTAINS} and \texttt{BELONGS\_TO}; the
sequential view sees only the three \texttt{NEXT\_*} edges; and so on.
This gives one node embedding per view. A node-conditioned attention
fusion then combines the five view-specific node embeddings into a
single embedding per node. After that, mean-and-max pooling over nodes
produces the graph embedding that the prediction head consumes.

\paragraph{Why this is called multiview.}
The term comes from multi-view learning, where the same input is
encoded through multiple distinct ``views'' (different feature spaces,
different modalities, or different relation types) and the resulting
representations are combined. In our case the views differ in which
edge types are visible. The same node set is shared, but each view's
GNN sees a different adjacency.

\paragraph{What inspired it.}
Two things. First, Joern's CPG explicitly separates AST, CFG and DDG
edges and uses them differently in queries. We wanted the same
separation at the model level: a single shared GNN that mixes all
edge types together would force syntax, sequence and coreference
through the same parameters, even though they describe quite different
relationships. Second, in the SemVul project (the predecessor of this
work) the multiview decomposition gave a measurable improvement on the
CPG side. We carried the same architecture across to the TPG, with
the views replaced by their text-domain analogues.

\paragraph{Why it is the right shape for the TPG.}
Adding security edges expands the edge vocabulary from 13 to 23 types.
With a single-view GNN, those 10 new types would have to share message
passing parameters with everything else. With the multiview design, the
security view is a separate encoder; if it has zero edges (because
\texttt{--include-security-edges} is off) the encoder is the identity
and contributes nothing, and if it has edges it learns its own
parameters without disturbing the other four views. That makes the
overlay clean and reversible.

\section{spaCy: How the Linguistic Backend Works and What We Take From It}

The TPG generic layer is built by a spaCy frontend. spaCy is what
turns raw text into the initial \textsc{document}, \textsc{paragraph},
\textsc{sentence}, \textsc{token}, \textsc{entity}, \textsc{predicate},
\textsc{argument}, \textsc{noun\_phrase}, \textsc{verb\_phrase} and
\textsc{clause} nodes, and the linguistic edges between them.

\subsection{What spaCy Is and How It Processes Text}

spaCy describes itself as ``Industrial-strength Natural Language
Processing in Python''.\footnote{\url{https://spacy.io/}} It ships
pre-trained pipelines for English (and many other languages) that
perform tokenisation, part-of-speech tagging, dependency parsing,
lemmatisation, named entity recognition, sentence segmentation, and
noun-phrase chunking, all in one pass.

The processing flow is:

\begin{enumerate}
    \item Raw text comes in as a string.
    \item The tokenizer splits it into tokens. spaCy's tokenizer is
    non-destructive: the tokenised \texttt{Doc} can reconstruct the
    original input character for character. This matters for the TPG
    because we anchor every node to character offsets, and we need
    those offsets to line up with the SecBERT embedder.
    \item Statistical components annotate each token: POS tag, lemma,
    dependency-parse head, dependency relation label, NER tag, IOB tag.
    \item Sentence boundaries are emitted as \texttt{Doc.sents}.
    \item Noun-phrase chunks are emitted as \texttt{Doc.noun\_chunks}
    when a parser is present.
\end{enumerate}

The frontend loads the English pipeline \texttt{en\_core\_web\_sm}
when available. If that model is not installed it falls back to a
blank English pipeline plus a sentencizer. The fallback can split
sentences but does not provide the parser, the lemmatiser, or the
NER. A flag is checked before every step that needs those components,
so the TPG degrades gracefully.

\subsection{What Counts as ``NLP'' Here}

When we say spaCy gives us NLP, we mean it provides four concrete
linguistic annotations on top of the raw token sequence:

\begin{itemize}
    \item \textbf{Tokenisation and POS tagging.} Each token gets a
    surface form, a lemma, and a part-of-speech tag.
    \item \textbf{Dependency parsing.} Each token is linked to a
    syntactic head, with a labelled relation
    (\texttt{nsubj}, \texttt{dobj}, \texttt{amod}, \texttt{prep}, etc.).
    This is what gives the TPG its \texttt{DEP} edges.
    \item \textbf{Named entity recognition.} Spans of tokens are
    classified as PERSON, ORG, GPE, PRODUCT, DATE, etc. These spans
    become \textsc{entity} nodes in the TPG.
    \item \textbf{Sentence and noun-phrase segmentation.} Sentences
    become \textsc{sentence} nodes; noun chunks become
    \textsc{noun\_phrase} nodes, attached to their constituent tokens
    via \texttt{BELONGS\_TO} edges.
\end{itemize}

These are statistical annotations produced by the trained pipeline.
The TPG then organises them into typed graph structure. spaCy
contributes the per-token signal; the TPG contributes the graph shape.

\subsection{The Exact Features We Take From spaCy}

The frontend extracts the following features and turns them into
nodes, properties or edges.

\begin{table}[h]
\centering
\small
\begin{tabular}{p{0.30\textwidth}p{0.30\textwidth}p{0.32\textwidth}}
\toprule
spaCy annotation & Used for & TPG output \\
\midrule
\texttt{Doc.sents} (sentence boundaries) & Splitting paragraphs into sentences & \textsc{sentence} nodes; \texttt{CONTAINS} edges from paragraph to sentence \\
Token text and \texttt{token.idx} & Per-token surface and character offset & \textsc{token} nodes; \texttt{char\_start} / \texttt{char\_end} properties \\
\texttt{token.lemma\_} & Canonical form of each token & \texttt{lemma} property on \textsc{token} \\
\texttt{token.pos\_} & Part-of-speech tag & \texttt{pos\_tag} property on \textsc{token} \\
\texttt{token.dep\_} & Dependency relation label (\texttt{nsubj}, \texttt{dobj}, etc.) & \texttt{dep\_rel} property on \textsc{token}; label on \texttt{DEP} edge \\
\texttt{token.head} (parser output) & Syntactic head of each token & \texttt{DEP} edges from head token to dependent \\
\texttt{token.ent\_type\_}, \texttt{token.ent\_iob\_} & NER tag and IOB position & properties on \textsc{token} for downstream entity span construction \\
\texttt{sent.ents} & Named entity spans & \textsc{entity} nodes; \texttt{CONTAINS} from sentence; \texttt{BELONGS\_TO} from entity to constituent tokens \\
Content verbs (token.pos\_ in VERB / AUX filter) & Predicate-argument extraction & \textsc{predicate} nodes; \texttt{CONTAINS} from sentence; \texttt{BELONGS\_TO} to verb token \\
Predicate children with mapped dependency labels (\texttt{nsubj}, \texttt{dobj}, etc.) & Semantic role assignment & \textsc{argument} nodes; \texttt{SRL\_ARG} edges from token and from argument node to predicate \\
\texttt{sent.noun\_chunks} & Noun-phrase chunking & \textsc{noun\_phrase} nodes; \texttt{CONTAINS} from sentence; \texttt{BELONGS\_TO} from chunk to its tokens \\
Clausal dependency labels (\texttt{advcl}, \texttt{relcl}, \texttt{ccomp}, \texttt{acl}) & Clause detection & \textsc{clause} nodes \\
Verb-phrase dependents (auxiliaries, negation, particles, adverbial modifiers) & Verb-phrase chunking & \textsc{verb\_phrase} nodes \\
Token order within a sentence & Linear flow & \texttt{NEXT\_TOKEN} edges \\
Sentence order within a paragraph & Linear flow & \texttt{NEXT\_SENT} edges \\
Paragraph order within the document & Linear flow & \texttt{NEXT\_PARA} edges \\
Last token of sentence $N$ to first token of sentence $N{+}1$ & Cross-sentence continuity & extra \texttt{NEXT\_TOKEN} edges \\
\bottomrule
\end{tabular}
\caption{spaCy annotations used by the TPG generic frontend, and what
the frontend produces from each one.}
\end{table}

After this, the security frontend overlays domain entities (CVE IDs,
CWE IDs, software, versions, attack vectors, impacts, severity,
remediations) on top of the spaCy graph. Then the SecBERT
embedder runs over the raw text, aligns its WordPiece sub-tokens to
the spaCy tokens via character offsets, and attaches a 768-dim
contextual vector to every \textsc{token} node (with mean-pooled
versions on \textsc{sentence} and \textsc{entity}).

\subsection{Why spaCy Specifically}

spaCy gives us the four things we need in one consistent
\texttt{Doc} object: tokenisation, POS and dependency parsing, NER,
and sentence segmentation. Other libraries provide some of these
(NLTK has tokenisation and POS; Stanza has dependency parsing);
spaCy provides all of them with reasonable accuracy and a stable API,
and the character offsets line up exactly with the BERT embedder we
use later. The TPG schema is decoupled from spaCy (the frontend
class is one of several possible frontends), so a different parser
could be substituted, but the mapping from spaCy outputs to TPG
structure is the one that the project actually implements and runs.

\section{How Security Entities and \texttt{SEC\_*} Edges Are Extracted}

Once spaCy has built the linguistic skeleton, the security frontend
runs a second pass over the same text and adds domain entities. The
extraction is deliberately rule-based: regular expressions for
patterns with stable surface forms, keyword dictionaries with
longest-first matching for everything else. Rules are precise on
canonical advisories and easy to audit. For surface forms that the
rules miss, a SecBERT-similarity step in the hybrid frontend adds
candidate entities, with rule entities winning on overlap.

\subsection{The Regular Expressions}

The patterns below are shown literally (in monospace, without LaTeX
interpretation). They are case-insensitive unless noted.

\paragraph{CVE identifier.}
\begin{verbatim}
CVE[<six Unicode hyphens>\-]\d{4}[<six Unicode hyphens>\-]\d{4,}
\end{verbatim}
Matches the canonical ASCII form (CVE followed by a hyphen, year,
hyphen, four-or-more digits) and the same form written with non-ASCII
hyphens. The character class contains all six Unicode hyphen code
points in the range U+2010 to U+2015 along with the ASCII hyphen,
because LLM-generated summaries often emit a non-ASCII hyphen between
the year and the sequence number.

\paragraph{CWE identifier.}
\begin{verbatim}
CWE-\d{1,4}
\end{verbatim}
Plain ASCII hyphen, weakness identifier of one to four digits.

\paragraph{Software version (semantic-version shape).}
\begin{verbatim}
\b\d+\.\d+(?:\.\d+)*(?:-[a-zA-Z0-9]+)?\b
\end{verbatim}
Catches \texttt{1.2.3}, \texttt{1.0.0-rc1}, \texttt{4.18.0-372}, etc.
Anchored on word boundaries. A separate context filter rejects numerals
that look like CVSS scores when the surrounding sentence mentions
``cvss'' or ``score''.

\paragraph{CVSS score in context.}
\begin{verbatim}
\b(?:CVSS[:\s]*v?\d?\.?\d?[:\s]*[/\-]?\s*)?(\d+\.\d+)(?:\s*/\s*10)?\b
\end{verbatim}
A composite pattern that matches \texttt{CVSS 9.8}, \texttt{CVSS:9.8},
\texttt{CVSS v3 9.8}, \texttt{9.8/10}, and parenthetical scores near
the word ``CVSS''.

\paragraph{Severity with context.}
\begin{verbatim}
\b(critical|high|medium|moderate|low|informational)\s+
   (severity|risk|impact|cvss\s*score|cvss|score|
    vulnerability|flaw|priority|confidence)\b
\end{verbatim}
Requires the severity adjective to be followed by a security-context
noun. The required second token is what stops phrases like ``high
availability'' or ``low memory'' from being tagged as severities.

\paragraph{Bare severity (fallback).}
\begin{verbatim}
\b(critical|high|medium|moderate|low)\b(?=\s*([.,;:!?\)\]]|$))
\end{verbatim}
A looser fallback that fires only when the severity adjective is
immediately followed by punctuation or end-of-clause. The extractor
runs this fallback only inside clauses that already contain a security
keyword within $\pm 60$ characters.

\paragraph{Code element.}
\begin{verbatim}
\b(strcpy|strcat|sprintf|gets|scanf|memcpy|memmove|
   malloc|free|printf|fprintf|eval|exec|system|popen|
   fork|[a-z_][a-z0-9_]*\(\))\b
\end{verbatim}
A list of dangerous C-stdlib functions and scripting eval/exec/system
names, plus the generic \texttt{name()} shape that catches any other
identifier written with parentheses.

\paragraph{Remediation actions (about fifteen patterns).}
The remediation extractor compiles a list of small patterns and tags
each match with a remediation type. The most common ones are:
\begin{verbatim}
upgrad\w*\s+to\s+(?:version\s+)?(\S+)         → upgrade
update\s+to\s+(?:version\s+)?(\S+)            → upgrade
patch\w*\s+(?:by\s+)?(?:applying|installing)  → patch
apply\s+(?:the\s+)?patch                      → patch
security\s+patch                              → patch
patched\s+version                             → patch
install(ing)?\s+the\s+update                  → update
\bfix(?:ed)?\s+in\s+(?:version\s+)?(\S+)      → update
workaround                                    → workaround
mitigat\w+                                    → mitigation
remediat\w+                                   → mitigation
disabl\w+\s+\w+                               → disable_feature
restrict\w*\s+\w+                             → restrict_access
firewall\s+(?:rule|block)                     → network_block
\end{verbatim}

\paragraph{Vendor and product fallback.}
A pattern that catches title-case ``Vendor Product'' mentions near a
version number. Used to recover niche products that are not in the
curated software dictionary (for example, proprietary appliances and
in-house tooling).

\subsection{The Dictionary-Based Extractors}

Some entities cannot be captured cleanly with a regex because their
surface forms are open-ended. For these we use curated keyword
dictionaries with longest-first matching:

\begin{itemize}
    \item \textbf{Software products.} A set of about 140 known names
    (Apache, nginx, OpenSSL, Linux kernel, WordPress, Spring Boot,
    Telerik, Confluence, vCenter, FortiOS, Palo Alto PAN-OS, etc.).
    Multi-word names like \texttt{apache http server} are matched
    before \texttt{apache} so the longer span wins. Word boundaries
    are anchored on either side so \texttt{git} does not match
    \texttt{github}.
    \item \textbf{Code constructs.} A focused list of around 35
    security-relevant phrases (\texttt{deserialization},
    \texttt{template engine}, \texttt{xml parser},
    \texttt{session token}, \texttt{prepared statement},
    \texttt{authorisation check}, etc.) that appear in CVE
    descriptions but are rarely written with parentheses.
    \item \textbf{Attack vectors.} About ten keyword phrases
    (\texttt{remote}, \texttt{local}, \texttt{authenticated user},
    \texttt{user-supplied input}, etc.).
    \item \textbf{Impact types.} About 25 phrases that name what an
    attacker gains (\texttt{remote code execution},
    \texttt{denial of service}, \texttt{information disclosure},
    \texttt{privilege escalation}, etc.).
    \item \textbf{Vulnerability types.} About 25 classification
    phrases (\texttt{buffer overflow}, \texttt{cross-site scripting},
    \texttt{sql injection}, \texttt{path traversal}, \texttt{use after free},
    \texttt{race condition}, etc.). The vulnerability-type extractor
    runs first and shares its claimed spans with the impact and
    attack-vector extractors, so a phrase that appears in two
    dictionaries is tagged once with the most specific label.
\end{itemize}

For each dictionary the extractor sorts entries longest-first,
compiles a word-boundary-anchored regex per entry, and keeps a
\texttt{seen\_spans} list so two patterns cannot claim the same
character range twice. Every entity that survives is added as an
\textsc{entity} node in the graph, tagged with its
\texttt{domain\_type} property, and linked to its containing
\textsc{sentence} via a \texttt{CONTAINS} edge.

\subsection{From Entities to \texttt{SEC\_*} Edges}

The entities are the nodes. The relations between them are edges.
The security-relations pass walks the entity nodes and emits ten
typed \texttt{SEC\_*} edges:

\begin{table}[h]
\centering
\small
\begin{tabular}{lll}
\toprule
Edge type & Source role & Target role \\
\midrule
\texttt{SEC\_AFFECTS}        & CVE / vulnerability    & SOFTWARE \\
\texttt{SEC\_HAS\_VERSION}   & SOFTWARE               & VERSION \\
\texttt{SEC\_LOCATED\_IN}    & VULN\_TYPE             & CODE\_ELEMENT \\
\texttt{SEC\_CLASSIFIED\_AS} & CVE                    & CWE \\
\texttt{SEC\_EXPLOITED\_BY}  & CVE / vulnerability    & ATTACK\_VECTOR \\
\texttt{SEC\_CAUSES}         & ATTACK\_VECTOR         & IMPACT \\
\texttt{SEC\_MITIGATED\_BY}  & CVE / vulnerability    & REMEDIATION \\
\texttt{SEC\_USES\_FUNCTION} & CVE / vulnerability    & CODE\_ELEMENT \\
\texttt{SEC\_THREATENS}      & ATTACK\_VECTOR         & SOFTWARE \\
\texttt{SEC\_HAS\_SEVERITY}  & CVE                    & SEVERITY \\
\bottomrule
\end{tabular}
\caption{The ten \texttt{SEC\_*} edge types emitted by the security
relations pass.}
\end{table}

The pass is deliberately conservative. It mostly looks for source and
target entities inside the same sentence (using each entity's
sentence-index property), and only links them if both are present.
This keeps the pass cheap and avoids spurious cross-document edges.
For relations where the typical writing pattern crosses one sentence
boundary (mitigation phrases at the end of an advisory, for example)
the pass widens the window slightly. The result is a small,
high-precision set of typed edges that the multiview GNN's security
view can consume directly.

\subsection{Why Rule-Based Extraction Is Used Here}

Two reasons. First, the surface forms that matter (CVE-IDs, CWE-IDs,
semantic versions, CVSS scores) are syntactically rigid. Regular
expressions are the natural fit and produce near-perfect precision on
canonical advisories. Second, the extracted entities feed an
interpretable downstream pipeline: a graph-level decision can be
traced back to specific entity nodes, and an analyst can audit the
extraction by looking at the regex that produced each node. A
neural extractor would make the same decisions opaque.

The trade-off is recall on novel surface forms. The hybrid frontend
addresses this by adding a SecBERT-similarity pass that recovers
entities the rules miss, with rule entities winning on overlap so the
audit trail stays intact.

\section{Training Hyperparameters and Configuration}

The training command used for the 20-run text-source ablation and the
NVD/KEV runs is:

\begin{center}
\texttt{python -m epss.run\_pipeline --backbone multiview --hybrid --label-mode soft --epochs 100 --no-epss-feature \dots}
\end{center}

The \texttt{\dots} is the variant-specific flag (\texttt{--summary-source},
\texttt{--summary-only-tpg}, \texttt{--include-summary-in-tpg}, etc.).
The defaults for all other hyperparameters are:

\begin{table}[h]
\centering
\small
\begin{tabular}{lll}
\toprule
Hyperparameter & CLI flag & Default \\
\midrule
GNN backbone               & \texttt{--backbone}        & \texttt{multiview} (we used this; alternatives: \texttt{gcn}, \texttt{gat}, \texttt{sage}, \texttt{edge\_type}, \texttt{rgat}) \\
Hidden dimension           & \texttt{--hidden}          & 128 \\
Number of GNN layers       & \texttt{--layers}          & 3 (the multiview backbone defaults to 4 GGNN unrolling steps internally) \\
Attention heads (GAT/RGAT) & \texttt{--heads}           & 4 \\
Dropout                    & \texttt{--dropout}         & 0.3 \\
Epochs                     & \texttt{--epochs}          & 100 (we set this explicitly) \\
Batch size                 & \texttt{--batch-size}      & 32 \\
Learning rate              & \texttt{--lr}              & $1 \times 10^{-3}$ \\
Weight decay               & \texttt{--weight-decay}    & $1 \times 10^{-4}$ \\
Early-stopping patience    & \texttt{--patience}        & 15 epochs without val PR-AUC improvement \\
Random seed                & \texttt{--seed}            & 42 \\
SecBERT embedding dim      & \texttt{--embedding-dim}   & 768 \\
Top-K CWEs encoded         & \texttt{--top-k-cwes}      & 25 \\
\bottomrule
\end{tabular}
\caption{Hyperparameters and their defaults.}
\end{table}

The training loop has the following pieces:

\begin{itemize}
    \item \textbf{Loss.} Binary cross-entropy with logits. For the soft
    EPSS target, BCE works on continuous labels in $[0, 1]$ as well as
    on hard $\{0, 1\}$ labels. A positive class weight is set to the
    negative-to-positive ratio in the training set so that the rare
    high-EPSS class gets enough gradient signal.
    \item \textbf{Optimiser.} AdamW with the learning rate and weight
    decay above.
    \item \textbf{Scheduler.} Reduce-on-plateau, mode max, factor 0.5,
    patience 5, minimum learning rate $10^{-6}$, monitoring validation
    PR-AUC.
    \item \textbf{Gradient clipping.} Gradient-norm clipping at 1.0
    after each backward pass.
    \item \textbf{Train loader.} Shuffle on, drop-last on. Dropping
    the last partial mini-batch avoids a known crash when the final
    batch has only one example and the model contains a BatchNorm
    layer in the prediction head (BN needs more than one element per
    channel).
    \item \textbf{Class weights.} Computed once from the full dataset
    before the splits are created. This is approximately the train-set
    ratio because the splits are stratified.
    \item \textbf{Splits.} If no external test file is provided, the
    dataset is split into train / val / test by class-stratified random
    sampling. If an external test file is provided, the internal split
    is train (85\%) / val (15\%) and the external file becomes the test
    set.
\end{itemize}

For the leakage-free configuration used in the 20-run ablation, the
relevant input dimensions are:

\begin{table}[h]
\centering
\small
\begin{tabular}{ll}
\toprule
Quantity & Value \\
\midrule
Node feature dimension $d_{\text{in}}$ & 781 (13-dim node-type one-hot $+$ 768-dim SecBERT) \\
Tabular feature dimension              & 55 (with EPSS removed) \\
Number of edge types                   & 13 (without security edges) or 23 (with) \\
Hidden width $h$                       & 128 \\
Tabular MLP width $h_t$                & 64 \\
Pooled graph embedding                 & $2h = 256$ (mean concatenated with max) \\
Fusion vector $z = [g; t]$             & $2h + h_t = 320$ \\
\bottomrule
\end{tabular}
\caption{Tensor dimensions in the leakage-free configuration.}
\end{table}

The companion document on the multiview formula writes out the
encoder, the node-conditioned attention fusion, and the hybrid
prediction head with all dimensions.

\section{Limitations}

Here are a few limitations:

\begin{itemize}
    \item Coreference is heuristic in the current enrichment pass. It
    uses recency and basic agreement, not a full neural coreference
    resolver.
    \item Discourse relations are marker-based. They catch explicit
    markers (``however'', ``because'', ``subsequently'') but do not
    perform full discourse parsing.
    \item Security entity extraction depends on regular expressions,
    keyword dictionaries, and SecBERT-similarity thresholds. New
    surface forms (Telegram threat-intelligence chatter, for example)
    are recovered less reliably than NVD-style advisories.
    \item Security edges are opt-in. If the pipeline does not enable
    them and the model's edge-type vocabulary is not expanded to 23,
    the GNN cannot use \texttt{SEC\_*} edges as first-class types.
    \item TPG is a representation. It improves the structural
    information available to the model. It does not guarantee a better
    PR-AUC on every downstream task. The 20-run ablation showed this
    plainly: the LLM summaries did not beat description-only on this
    in-distribution corpus.
\end{itemize}

\section{Conclusion}

The Text Property Graph is a structured, typed, attributed graph
representation of text. The central idea is borrowed from the Code
Property Graph: keep several views of the input in one graph rather
than as separate analyses. For text those views are dependency
structure, sequence, coreference, predicate-argument structure,
discourse, entity relation, and optional security relations.

We chose a GNN over the TPG because the input is heterogeneous and
relational. We chose the TPG over a single-view alternative because
the CPG inspiration generalises naturally to text. We chose the
multiview backbone because different edge groups deserve different
encoders. We chose spaCy as the linguistic backend because it gives
us tokenisation, parsing, NER and sentence segmentation in one
consistent \texttt{Doc} object whose character offsets line up with
the SecBERT embedder.

For security text, this combination preserves the relations that
matter (which CVE affects which software, with which attack vector,
causing which impact, mitigated by which fix) and feeds them to the
model as typed graph structure rather than as flat embeddings. The
empirical results show what this gives us in practice: a strong
description-driven baseline whose signal does not get displaced by
adding more text, and a representation that is ready for the
cross-modal and temporal experiments planned next.

\end{document}
